\renewcommand{\thesection}{\arabic{section}}\renewcommand{\theHsection}{chapter.\arabic{section}}\setcounter{section}{4}\section{单变量函数的积分学}
\noindent\begin{minipage}{\textwidth}\ProblemHeading{例题}{5.1（改写）}
设 $a<b$，$c$ 为常数。用 Riemann 积分的定义证明 $\displaystyle\int_a^b c\,dx=c(b-a)$。
\end{minipage}\par\vfill
\noindent\begin{minipage}{\textwidth}\ProblemHeading{例题}{5.2}
设 $a<b$。用积分定义证明 $\displaystyle\int_a^b x\,dx=\frac12(b^2-a^2)$。
\end{minipage}\par\vfill
\noindent\begin{minipage}{\textwidth}\ProblemHeading{例题}{5.3}
设 $a<b$。用积分定义证明 $\displaystyle\int_a^b x^2\,dx=\frac13(b^3-a^3)$。
\end{minipage}\par\vfill
\noindent\begin{minipage}{\textwidth}\ProblemHeading{例题}{5.4}
设 $a<b$。用积分定义求 $\displaystyle\int_a^b\sin x\,dx$。
\end{minipage}\par\vfill
\newpage
\noindent\begin{minipage}{\textwidth}\ProblemHeading{例题}{5.5}
设 $a<b$，$a\le c\le b$，函数\[J(x)=\begin{cases}0,&x\in[a,b],\ x\ne c,\\1,&x=c.\end{cases}\]用积分定义证明 $\displaystyle\int_a^bJ(x)\,dx=0$。
\end{minipage}\par\vfill
\noindent\begin{minipage}{\textwidth}\ProblemHeading{习题}{5.6}
设 $f,g$ 在 $[a,b]$ 上 Riemann 可积，$\alpha,\beta$ 为常数。证明 $\alpha f+\beta g$ 也可积，且\[\int_a^b(\alpha f(x)+\beta g(x))\,dx=\alpha\int_a^b f(x)\,dx+\beta\int_a^b g(x)\,dx.\]
\end{minipage}\par\vfill
\noindent\begin{minipage}{\textwidth}\ProblemHeading{习题}{5.7}
设 $a<b$。自行绘图，解释下列积分性质的几何意义。
\SubQ{1}{若 $f,g$ 在 $[a,b]$ 上可积，则 $f\ge0$ 时 $\int_a^b f\ge0$；$f\le g$ 时 $\int_a^b f\le\int_a^b g$；并有 $|\int_a^b f|\le\int_a^b|f|$。}
\SubQ{2}{若 $f$ 在 $[a,b]$ 上可积，且 $m\le f(x)\le M$，则 $m(b-a)\le\int_a^b f(x)\,dx\le M(b-a)$。}
\SubQ{3}{若 $f$ 在 $[a,b]$ 上连续，则存在 $\xi\in[a,b]$，使 $\int_a^b f(x)\,dx=f(\xi)(b-a)$。}
\end{minipage}\par\vfill
\noindent\begin{minipage}{\textwidth}\ProblemHeading{例题}{5.8}
计算 $\displaystyle\int_0^2|x^3-1|\,dx$。
\end{minipage}\par\vfill
\newpage
\noindent\begin{minipage}{\textwidth}\ProblemHeading{例题}{5.9}
计算 $\displaystyle\int_{-1}^1\operatorname{sgn}x\,dx$，其中 $\operatorname{sgn}x$ 在 $x>0,x=0,x<0$ 时分别为 $1,0,-1$。
\end{minipage}\par\vfill
\noindent\begin{minipage}{\textwidth}\ProblemHeading{例题}{5.10}
计算 $\displaystyle\int_{-1}^0\frac{x^2+1}{x^4+1}\,dx$。
\end{minipage}\par\vfill
\noindent\begin{minipage}{\textwidth}\ProblemHeading{例题}{5.11}
设 $a>0$，求 $\displaystyle\int_0^a\sqrt{a^2-x^2}\,dx$。
\end{minipage}\par\vfill
\noindent\begin{minipage}{\textwidth}\ProblemHeading{例题}{5.12}
计算 $\displaystyle\int_0^{\pi/2}\sin^n x\,dx$ 及 $\displaystyle\int_0^{\pi/2}\cos^n x\,dx$，其中 $n=0,1,2,\ldots$。
\end{minipage}\par\vfill
\newpage
\noindent\begin{minipage}{\textwidth}\ProblemHeading{例题}{5.13}
计算 $\displaystyle\int_0^1\frac{\arctan x}{1+x}\,dx$。
\end{minipage}\par\vfill
\newpage
\noindent\begin{minipage}{\textwidth}\ProblemHeading{习题}{5.14}
判断下列函数在 $[0,1]$ 上是否 Riemann 可积，并说明理由，其中 $[t]$ 为不大于 $t$ 的最大整数。
\SubQ{1}{$f(x)=\begin{cases}\sin x/x,&x\ne0,\\1,&x=0.\end{cases}$}
\SubQ{2}{$f(x)=\begin{cases}x^{-1}\sin(1/x),&x\ne0,\\1,&x=0.\end{cases}$}
\SubQ{3}{$f(x)=\begin{cases}1/[1/x],&0<x\le1,\\0,&x=0.\end{cases}$}
\end{minipage}\par\vfill
\newpage
\noindent\begin{minipage}{\textwidth}\ProblemHeading{习题}{5.15}
设 $a<b$，$D(x)=1$（$x\in\mathbb Q$），$D(x)=0$（$x\notin\mathbb Q$）。证明 Dirichlet 函数 $D$ 在 $[a,b]$ 上不可积。
\end{minipage}\par\vfill
\noindent\begin{minipage}{\textwidth}\ProblemHeading{习题}{5.16}
举例说明：$|f|$ 在 $[a,b]$ 上可积，不能保证 $f$ 在该区间可积。
\end{minipage}\par\vfill
\noindent\begin{minipage}{\textwidth}\ProblemHeading{习题}{5.17}
设 $a<b$，证明或举例：
\SubQ{1}{若 $f$ 在 $[a,b]$ 上可积且非负，在某点 $c\in[a,b]$ 连续且 $f(c)>0$，则 $\int_a^b f(x)\,dx>0$。}
\SubQ{2}{若 $f$ 在 $[a,b]$ 上连续、非负且不恒为零，则 $\int_a^b f(x)\,dx>0$。}
\SubQ{3}{构造在 $[a,b]$ 上可积、非负且不恒为零，但积分为零的函数。}
\end{minipage}\par\vfill
\noindent\begin{minipage}{\textwidth}\ProblemHeading{习题}{5.18}
设 $a<b$，$f$ 在 $[a,b]$ 上单调递增。证明\[(b-a)f(a)\le\int_a^b f(x)\,dx\le(b-a)f(b).\]若 $f$ 单调递减，写出并证明相应不等式。
\end{minipage}\par\vfill
\newpage
\noindent\begin{minipage}{\textwidth}\ProblemHeading{习题}{5.19}
证明下列不等式：
\SubQ{1}{$\displaystyle\int_0^{2\pi}|a\sin x+b\cos x|\,dx\le2\pi\sqrt{a^2+b^2}$，其中 $a,b$ 为实常数。}
\SubQ{2}{$\displaystyle\int_0^1x^m(1-x)^n\,dx\le\frac{m^mn^n}{(m+n)^{m+n}}$，其中 $m,n>0$。}
\end{minipage}\par\vfill
\noindent\begin{minipage}{\textwidth}\ProblemHeading{习题}{5.20}
对积分中值定理，完成下列问题：
\SubQ{1}{若 $f\in C[a,b]$ 且 $a<b$，证明存在 $\xi\in(a,b)$，使 $\int_a^b f(x)\,dx=f(\xi)(b-a)$。}
\SubQ{2}{举例说明：若去掉连续性条件，仅假设可积，上述等式可能对所有 $\xi\in[a,b]$ 都不成立。}
\end{minipage}\par\vfill
\noindent\begin{minipage}{\textwidth}\ProblemHeading{习题}{5.21}
设 $a<b$，$f\in C[a,b]$ 且 $\int_a^b f(x)\,dx=0$。证明 $f$ 在 $(a,b)$ 内至少有一个零点。
\end{minipage}\par\vfill
\noindent\begin{minipage}{\textwidth}\ProblemHeading{习题}{5.22}
设 $a<b$，完成下列问题：
\SubQ{1}{若 $f\in C[a,b]$，$g$ 在 $[a,b]$ 上可积且非负（或非正），证明存在 $\xi\in[a,b]$，使 $\int_a^b f(x)g(x)\,dx=f(\xi)\int_a^b g(x)\,dx$。}
\SubQ{2}{举例说明：若允许可积函数 $g$ 改变符号，上述结论未必成立。}
\end{minipage}\par\vfill
\newpage
\noindent\begin{minipage}{\textwidth}\ProblemHeading{习题}{5.23}
举例说明：可积函数 $f$ 在 $c$ 处连续，并不是 $F(x)=\int_a^x f(t)\,dt$ 在 $c$ 处可导的必要条件。
\end{minipage}\par\vfill
\noindent\begin{minipage}{\textwidth}\ProblemHeading{习题}{5.24}
求下列函数的导数：
\SubQ{1}{$f(x)=\displaystyle\int_0^{x^2}\sin(t^2)\,dt$。}
\SubQ{2}{$f(x)=\displaystyle\int_x^1\frac{dt}{1+t^2+\cos^2t}$。}
\SubQ{3}{$f(x)=\displaystyle\int_x^{x^2}e^{-t^2}\,dt$。}
\end{minipage}\par\vfill
\noindent\begin{minipage}{\textwidth}\ProblemHeading{习题}{5.25}
求下列函数的导数：
\SubQ{1}{$f(x)=\sin\left(\displaystyle\int_0^x\sin\left(\int_0^y\sin(t^2)\,dt\right)dy\right)$。}
\end{minipage}\par\vfill
\noindent\begin{minipage}{\textwidth}\ProblemHeading{习题}{5.26}
对下列函数求 $(f^{-1})^{\prime}(0)$，其中 $f^{-1}$ 为反函数：
\SubQ{1}{$f(x)=\displaystyle\int_0^x(1+\sin(\sin t))\,dt$。}
\SubQ{2}{$f(x)=\displaystyle\int_1^x e^{-t^2}\,dt$。}
\end{minipage}\par\vfill
\newpage
\noindent\begin{minipage}{\textwidth}\ProblemHeading{习题}{5.27}
设 $f$ 处处连续，$F(x)=\int_0^x x f(t)\,dt$。求 $F^{\prime}(x)$。
\end{minipage}\par\vfill
\noindent\begin{minipage}{\textwidth}\ProblemHeading{习题}{5.28}
设 $f$ 在 $[0,+\infty)$ 上连续且为正值。证明\[G(x)=\frac{\int_0^x t f(t)\,dt}{\int_0^x f(t)\,dt}\]在 $x>0$ 时单调递增。
\end{minipage}\par\vfill
\noindent\begin{minipage}{\textwidth}\ProblemHeading{习题}{5.29}
用 Newton-Leibniz 公式计算下列积分：
\SubQ{1}{$\displaystyle\int_0^\pi\sin x\,dx$。}
\SubQ{2}{$\displaystyle\int_0^1x^\alpha\,dx$，其中 $\alpha>0$。}
\SubQ{3}{$\displaystyle\int_1^2\ln x\,dx$。}
\end{minipage}\par\vfill
\noindent\begin{minipage}{\textwidth}\ProblemHeading{习题}{5.30}
用 Newton-Leibniz 公式计算下列积分：
\SubQ{1}{$\displaystyle\int_2^3\frac{dx}{2x^2+3x-2}$。}
\end{minipage}\par\vfill
\newpage
\noindent\begin{minipage}{\textwidth}\ProblemHeading{习题}{5.31}
设 $f(x)=\begin{cases}-1,&-1\le x<0,\\0,&x=0,\\1,&0<x\le1.\end{cases}$ 求 $F(x)=\int_{-1}^x f(t)\,dt$，并研究其在 $[-1,1]$ 上的可微性（端点取单侧导数）。
\end{minipage}\par\vfill
\noindent\begin{minipage}{\textwidth}\ProblemHeading{习题}{5.32}
计算下列平面图形的面积：
\SubQ{1}{曲线 $y=\sqrt x$ 与 $y=x^2$ 围成的图形。}
\SubQ{2}{曲线 $y=x^2$、$y=x^2/4$ 及直线 $y=1$ 围成的图形。}
\end{minipage}\par\vfill
\noindent\begin{minipage}{\textwidth}\ProblemHeading{习题}{5.33}
求下列极限：
\SubQ{1}{$\displaystyle\lim_{x\to0}\frac{\int_0^x\sin(t^3)\,dt}{x^4}$。}
\SubQ{2}{$\displaystyle\lim_{x\to0}\frac1{\sin^3x}\int_0^{\tan x}\arcsin(t^2)\,dt$。}
\SubQ{3}{$\displaystyle\lim_{n\to\infty}\sum_{k=0}^{n-1}\frac1{\sqrt{n^2-k^2}}$。}
\end{minipage}\par\vfill
\noindent\begin{minipage}{\textwidth}\ProblemHeading{习题}{5.34}
求下列极限：
\SubQ{1}{$\displaystyle\lim_{n\to\infty}\frac{1^p+2^p+\cdots+n^p}{n^{p+1}}$，其中 $p>0$。}
\end{minipage}\par\vfill
\newpage
\noindent\begin{minipage}{\textwidth}\ProblemHeading{习题}{5.35}
求下列极限：
\SubQ{1}{$\displaystyle\lim_{n\to\infty}\int_a^b e^{-nx^2}\,dx$，其中 $0<a<b$。}
\SubQ{2}{$\displaystyle\lim_{n\to\infty}\int_0^1\frac{x^n}{1+x}\,dx$。}
\SubQ{3}{$\displaystyle\lim_{n\to\infty}\int_n^{n+a}\frac{\sin x}{x}\,dx$，其中 $a>0$。}
\end{minipage}\par\vfill
\noindent\begin{minipage}{\textwidth}\ProblemHeading{习题}{5.36}
设 $a>0$，$f\in C[-a,a]$。证明：
\SubQ{1}{若 $f$ 为奇函数，则 $\int_{-a}^a f(x)\,dx=0$。}
\SubQ{2}{若 $f$ 为偶函数，则 $\int_{-a}^a f(x)\,dx=2\int_0^a f(x)\,dx$。}
\end{minipage}\par\vfill
\noindent\begin{minipage}{\textwidth}\ProblemHeading{习题}{5.37}
设连续函数 $f$ 以 $T>0$ 为周期。证明对任意实数 $a$，有\[\int_a^{a+T}f(x)\,dx=\int_0^T f(x)\,dx.\]
\end{minipage}\par\vfill
\noindent\begin{minipage}{\textwidth}\ProblemHeading{习题}{5.38}
计算下列积分。其中 $[x]$ 表示取整，出现的参数 $a,b>0$；遇无定义的端点或内部点时按连续延拓或反常积分解释。
\SubQ{1}{$\displaystyle\int_0^{2\pi}|\cos x|\,dx$。}
\SubQ{2}{$\displaystyle\int_{-3}^4[x]\,dx$。}
\SubQ{3}{$\displaystyle\int_{-1}^1\cos x\ln\frac{1+x}{1-x}\,dx$。}
\end{minipage}\par\vfill
\newpage
\noindent\begin{minipage}{\textwidth}\ProblemHeading{习题}{5.39}
计算下列积分。其中 $[x]$ 表示取整，出现的参数 $a,b>0$；遇无定义的端点或内部点时按连续延拓或反常积分解释。
\SubQ{1}{$\displaystyle\int_{-\pi/2}^{\pi/2}\frac{\cos^3x}{1+e^x}\,dx$。}
\SubQ{2}{$\displaystyle\int_0^{\ln2}\sqrt{1-e^{-2x}}\,dx$。}
\SubQ{3}{$\displaystyle\int_0^1x\arcsin x\,dx$。}
\end{minipage}\par\vfill
\noindent\begin{minipage}{\textwidth}\ProblemHeading{习题}{5.40}
计算下列积分。其中 $[x]$ 表示取整，出现的参数 $a,b>0$；遇无定义的端点或内部点时按连续延拓或反常积分解释。
\SubQ{1}{$\displaystyle\int_0^1x^3e^x\,dx$。}
\SubQ{2}{$\displaystyle\int_0^a\frac{dx}{x+\sqrt{a^2-x^2}}$。}
\SubQ{3}{$\displaystyle\int_0^{\pi/4}\sqrt{\tan x}\,dx$。}
\end{minipage}\par\vfill
\noindent\begin{minipage}{\textwidth}\ProblemHeading{习题}{5.41}
计算下列积分。其中 $[x]$ 表示取整，出现的参数 $a,b>0$；遇无定义的端点或内部点时按连续延拓或反常积分解释。
\SubQ{1}{$\displaystyle\int_0^{\pi/2}\frac{dx}{a^2\sin^2x+b^2\cos^2x}$。}
\SubQ{2}{$\displaystyle\int_{-1}^1x^4\sqrt{1-x^2}\,dx$。}
\SubQ{3}{$\displaystyle\int_0^{2\pi}\sin^6x\,dx$。}
\end{minipage}\par\vfill
\noindent\begin{minipage}{\textwidth}\ProblemHeading{习题}{5.42}
计算下列积分。其中 $[x]$ 表示取整，出现的参数 $a,b>0$；遇无定义的端点或内部点时按连续延拓或反常积分解释。
\SubQ{1}{$\displaystyle\int_{-1}^1e^{|x|}\arctan(e^x)\,dx$。}
\SubQ{2}{$\displaystyle\int_0^\pi\frac{\sec^2x}{2+\tan^2x}\,dx$。}
\end{minipage}\par\vfill
\newpage
\noindent\begin{minipage}{\textwidth}\ProblemHeading{习题}{5.43}
设 $f\in C[0,1]$。证明\[\int_0^\pi x f(\sin x)\,dx=\pi\int_0^{\pi/2}f(\sin x)\,dx,\]并计算 $\displaystyle\int_0^\pi\frac{x\sin x}{1+\cos^2x}\,dx$。
\end{minipage}\par\vfill
\noindent\begin{minipage}{\textwidth}\ProblemHeading{习题}{5.44}
证明 $\displaystyle\frac16<\int_0^1\sin(x^2)\,dx<\frac13$。
\end{minipage}\par\vfill
\noindent\begin{minipage}{\textwidth}\ProblemHeading{习题}{5.45}
函数 $f$ 在 $[a,b]$ 上的平均值定义为 $\frac1{b-a}\int_a^b f(x)\,dx$。分别求下列函数在 $[0,1]$ 及 $[0,10^5]$ 上的平均值、最大值和最小值。
\SubQ{1}{$f(x)=x$。}
\SubQ{2}{$f(x)=e^{-x}$。}
\SubQ{3}{$f(x)=xe^{-x}$。}
\end{minipage}\par\vfill
\noindent\begin{minipage}{\textwidth}\ProblemHeading{习题}{5.46}
考虑 $I=\displaystyle\int_0^{100}\frac{e^{-x}}{x+100}\,dx$。
\SubQ{1}{尽可能精确地估计 $I$ 的上、下界。}
\SubQ{2}{求 $I$ 的近似值，精确到 $0.0001$。}
\end{minipage}\par\vfill
\newpage
\noindent\begin{minipage}{\textwidth}\ProblemHeading{习题}{5.47}
对任意 $\alpha\in(0,1)$，证明以下两种条件下均有 $\int_0^\alpha f(x)\,dx\ge\alpha\int_0^1 f(x)\,dx$：
\SubQ{1}{$f$ 在 $[0,1]$ 上单调递减且连续。}
\SubQ{2}{仅假设 $f$ 在 $[0,1]$ 上单调递减。}
\end{minipage}\par\vfill
\noindent\begin{minipage}{\textwidth}\ProblemHeading{习题}{5.48}
设 $a<b$，$f$ 在 $[a,b]$ 上可微且 $|f^{\prime}(x)|\le M$。证明：
\SubQ{1}{若 $f(a)=0$，则 $\int_a^b|f(x)|\,dx\le\frac M2(b-a)^2$。}
\SubQ{2}{若 $f(a)=f(b)=0$，则 $\int_a^b|f(x)|\,dx\le\frac M4(b-a)^2$。}
\end{minipage}\par\vfill
\noindent\begin{minipage}{\textwidth}\ProblemHeading{习题}{5.49}
利用代换 $u=\sin x$，将 $\displaystyle\int\frac{dx}{\sqrt{\cos2x}}$ 化为椭圆积分的形式。
\end{minipage}\par\vfill
\noindent\begin{minipage}{\textwidth}\ProblemHeading{习题}{5.50}
设 $f$ 在包含 $a,x$ 的闭区间上具有 $n+1$ 阶连续导数。从\[f(x)-f(a)=\int_a^x f^{\prime}(t)\,dt=\int_a^x f^{\prime}(t)\,d(t-x)\]出发，反复分部积分，推导 $n$ 阶 Taylor 展开式及其余项的积分表示式。
\end{minipage}\par\vfill
\newpage
\noindent\begin{minipage}{\textwidth}\ProblemHeading{习题}{5.51}
设 $f\in C(\mathbb R)$，$g(x,y)=\int_0^x\bigl(f(t+y)-f(t)\bigr)\,dt$。证明 $g(x,y)=g(y,x)$。
\end{minipage}\par\vfill
\noindent\begin{minipage}{\textwidth}\ProblemHeading{例题}{5.52}
证明：闭区间 $[a,b]$ 上的连续函数一定 Riemann 可积。
\end{minipage}\par\vfill
\noindent\begin{minipage}{\textwidth}\ProblemHeading{例题}{5.53}
证明：在 $[a,b]$ 上有界且至多有有限个间断点的函数一定 Riemann 可积。
\end{minipage}\par\vfill
\noindent\begin{minipage}{\textwidth}\ProblemHeading{例题}{5.54}
证明：闭区间 $[a,b]$ 上的单调函数一定 Riemann 可积。
\end{minipage}\par\vfill
\newpage
\noindent\begin{minipage}{\textwidth}\ProblemHeading{例题}{5.55}
设 $a<c<b$。证明 $f$ 在 $[a,b]$ 上可积，当且仅当它分别在 $[a,c]$ 和 $[c,b]$ 上可积；并证明\[\int_a^b f(x)\,dx=\int_a^c f(x)\,dx+\int_c^b f(x)\,dx.\]
\end{minipage}\par\vfill
\noindent\begin{minipage}{\textwidth}\ProblemHeading{例题}{5.56}
证明：若 $f,g$ 在 $[a,b]$ 上 Riemann 可积，则乘积 $fg$ 也在 $[a,b]$ 上可积。
\end{minipage}\par\vfill
\noindent\begin{minipage}{\textwidth}\ProblemHeading{习题}{5.57}
设 $D(x)=1$（$x\in\mathbb Q$），$D(x)=0$（$x\notin\mathbb Q$）。求 $D$ 在 $[0,1]$ 上的 Darboux 上积分和下积分。
\end{minipage}\par\vfill
\noindent\begin{minipage}{\textwidth}\ProblemHeading{习题}{5.58}
给出 Darboux 上和与下和的几何解释。这里，对于 $[a,b]$ 的分割 $a=x_0<\cdots<x_n=b$，上和为 $\sum_i M_i(x_i-x_{i-1})$，下和为 $\sum_i m_i(x_i-x_{i-1})$，$M_i,m_i$ 分别为函数在第 $i$ 个小区间上的上、下确界。
\end{minipage}\par\vfill
\newpage
\noindent\begin{minipage}{\textwidth}\ProblemHeading{习题}{5.59}
证明：若 $f$ 在 $[a,b]$ 上可积，则 $|f|$ 也可积，并且 $\left|\int_a^b f(x)\,dx\right|\le\int_a^b|f(x)|\,dx$。
\end{minipage}\par\vfill
\noindent\begin{minipage}{\textwidth}\ProblemHeading{习题}{5.60}
设 $f$ 在 $[a,b]$ 上可积，且存在常数 $c>0$ 使 $f(x)\ge c$。证明 $1/f$ 在 $[a,b]$ 上可积。
\end{minipage}\par\vfill
\noindent\begin{minipage}{\textwidth}\ProblemHeading{例题}{5.61}
半径为 $R$ 米的半球形水池开口水平朝上，池内注满水。求将水全部抽到池沿所需做的功。水的密度为 $1\,\mathrm{t}/\mathrm m^3$，重力加速度记为 $g$。
\end{minipage}\par\vfill
\noindent\begin{minipage}{\textwidth}\ProblemHeading{例题}{5.62}
一根均匀细棒质量为 $m$、长度为 $2l$。单位质量的质点 $Q$ 位于棒所在直线上，距棒中心 $a>l$。求棒对质点 $Q$ 的引力大小。万有引力常量记为 $G$。
\end{minipage}\par\vfill
\newpage
\noindent\begin{minipage}{\textwidth}\ProblemHeading{习题}{5.63}
求下列曲线弧的长度，其中 $a>0$：
\SubQ{1}{抛物线 $y=x^2$ 从 $x=-a$ 到 $x=a$ 的弧。}
\SubQ{2}{星形线 $x=a\cos^3t$，$y=a\sin^3t$，$0\le t<2\pi$。}
\SubQ{3}{Archimedes 螺线 $r=a\theta$，$0\le\theta\le2\pi$。}
\end{minipage}\par\vfill
\noindent\begin{minipage}{\textwidth}\ProblemHeading{习题}{5.64}
计算下列曲线围成的平面图形的面积，其中 $a>0$：
\SubQ{1}{双纽线 $r^2=a^2\cos2\theta$，$\theta\in[-\pi/4,\pi/4]\cup[3\pi/4,5\pi/4]$。}
\SubQ{2}{曲线 $x=t-\sin t$，$y=1-\cos t$（$0\le t\le2\pi$）与 $x$ 轴。}
\SubQ{3}{曲线 $y=e^x$、$y=e^{-x}$ 及直线 $x=1$。}
\end{minipage}\par\vfill
\noindent\begin{minipage}{\textwidth}\ProblemHeading{习题}{5.65}
计算下列旋转体的体积：
\SubQ{1}{$y=\sin x$（$0\le x\le\pi$）与 $x$ 轴围成的图形，分别绕 $x$ 轴及 $y$ 轴旋转一周。}
\SubQ{2}{$y=e^{x^2}$（$0\le x\le1$）与 $x$ 轴围成的图形，绕 $y$ 轴旋转一周。}
\SubQ{3}{$x=t-\sin t$，$y=1-\cos t$（$0\le t\le2\pi$）与 $x$ 轴围成的图形，绕 $x$ 轴旋转一周。}
\end{minipage}\par\vfill
\noindent\begin{minipage}{\textwidth}\ProblemHeading{习题}{5.66}
证明：半径为 $R$ 的球中，高为 $h$ 的球缺体积为 $\pi h^2(R-h/3)$，其中 $0\le h\le2R$。
\end{minipage}\par\vfill
\newpage
\noindent\begin{minipage}{\textwidth}\ProblemHeading{习题}{5.67}
求下列曲线旋转一周所得旋转曲面的面积（重合的部分只计一次），其中 $r,a,b>0$：
\SubQ{1}{$x^2+y^2=r^2$ 绕 $x$ 轴。}
\SubQ{2}{$x^2/a^2+y^2/b^2=1$ 绕 $y$ 轴，其中 $a>b$。}
\SubQ{3}{$y=a\cosh x$（$0\le x\le a$）绕 $x$ 轴。}
\end{minipage}\par\vfill
\noindent\begin{minipage}{\textwidth}\ProblemHeading{习题}{5.68}
求下列曲线旋转一周所得旋转曲面的面积（重合的部分只计一次），其中 $r,a,b>0$：
\SubQ{1}{$r=a(1+\cos\theta)$ 绕极轴。}
\end{minipage}\par\vfill
\noindent\begin{minipage}{\textwidth}\ProblemHeading{习题}{5.69}
半径为 $r$ 的球完全浸入水中，球顶与水面相切；球的密度为1（水的密度同取为1）。缓慢竖直提升球，直至球底刚好到达水面，需做多少功？重力加速度记为 $g$。
\end{minipage}\par\vfill
\noindent\begin{minipage}{\textwidth}\ProblemHeading{习题}{5.70}
两根长度均为 $l$、质量均为 $m$ 的均匀细杆位于同一直线上，两杆最近端点的距离为 $l$。求两杆之间的引力大小，万有引力常量记为 $G$。
\end{minipage}\par\vfill
\newpage
\noindent\begin{minipage}{\textwidth}\ProblemHeading{例题}{5.71}
设 $\alpha>0$。证明 $\displaystyle\int_0^{+\infty}\frac{dx}{(1+x^2)(1+x^\alpha)}$ 收敛，并求其值。
\end{minipage}\par\vfill
\noindent\begin{minipage}{\textwidth}\ProblemHeading{例题}{5.72}
证明 $\displaystyle\int_0^{\pi/2}\ln\sin x\,dx$ 收敛，并求其值。
\end{minipage}\par\vfill
\noindent\begin{minipage}{\textwidth}\ProblemHeading{习题}{5.73}
判断下列反常积分是否收敛，若收敛则求其值。其中 $n$ 为非负整数。
\SubQ{1}{$\displaystyle\int_0^{+\infty}xe^{-x^2}\,dx$。}
\SubQ{2}{$\displaystyle\int_0^{+\infty}x\sin x\,dx$。}
\SubQ{3}{$\displaystyle\int_2^{+\infty}\frac{\ln x}{x}\,dx$。}
\end{minipage}\par\vfill
\noindent\begin{minipage}{\textwidth}\ProblemHeading{习题}{5.74}
判断下列反常积分是否收敛，若收敛则求其值。其中 $n$ 为非负整数。
\SubQ{1}{$\displaystyle\int_1^{+\infty}\frac{\arctan x}{x}\,dx$。}
\SubQ{2}{$\displaystyle\int_0^{+\infty}e^{-x}\sin x\,dx$。}
\SubQ{3}{$\displaystyle\int_{-\infty}^{+\infty}\frac{dx}{x^2+2x+2}$。}
\end{minipage}\par\vfill
\newpage
\noindent\begin{minipage}{\textwidth}\ProblemHeading{习题}{5.75}
判断下列反常积分是否收敛，若收敛则求其值。其中 $n$ 为非负整数。
\SubQ{1}{$\displaystyle\int_0^1\ln x\,dx$。}
\SubQ{2}{$\displaystyle\int_{-1}^1\frac{dx}{\sqrt{1-x^2}}$。}
\SubQ{3}{$\displaystyle\int_0^1\frac{x\ln x}{(1-x^2)^{3/2}}\,dx$。}
\end{minipage}\par\vfill
\noindent\begin{minipage}{\textwidth}\ProblemHeading{习题}{5.76}
判断下列反常积分是否收敛，若收敛则求其值。其中 $n$ 为非负整数。
\SubQ{1}{$\displaystyle\int_0^1\ln\frac1{1-x^2}\,dx$。}
\SubQ{2}{$\displaystyle\int_0^{+\infty}x^n e^{-x}\,dx$。}
\SubQ{3}{$\displaystyle\int_0^1(\ln x)^n\,dx$。}
\end{minipage}\par\vfill
\noindent\begin{minipage}{\textwidth}\ProblemHeading{习题}{5.77}
Cauchy 主值定义为 $\operatorname{PV}\int_{-\infty}^{+\infty}f(x)\,dx=\lim_{b\to+\infty}\int_{-b}^b f(x)\,dx$。研究下列主值是否收敛：
\SubQ{1}{$\operatorname{PV}\displaystyle\int_{-\infty}^{+\infty}\frac{x}{1+x^2}\,dx$。}
\SubQ{2}{$\operatorname{PV}\displaystyle\int_{-\infty}^{+\infty}\frac{|x|}{1+x^2}\,dx$。}
\end{minipage}\par\vfill
\noindent\begin{minipage}{\textwidth}\ProblemHeading{习题}{5.78}
反常积分 $\int_a^{+\infty}f$ 在 $a$ 处也可能有瑕点，此时按 $\int_a^b f+\int_b^{+\infty}f$ 定义，要求两项分别收敛（$b>a$）。
\SubQ{1}{证明 $\displaystyle\int_1^{+\infty}\frac{dx}{x\sqrt{x-1}}$ 收敛，并求其值。}
\SubQ{2}{证明对任意实数 $\alpha$，$\displaystyle\int_0^{+\infty}x^{-\alpha}\,dx$ 发散。}
\end{minipage}\par\vfill
\newpage
\noindent\begin{minipage}{\textwidth}\ProblemHeading{习题}{5.79}
若 $c\in(a,b)$ 是被积函数的内部瑕点，积分 $\int_a^b f$ 按 $\int_a^c f+\int_c^b f$ 定义，要求两项分别收敛。判断下列积分是否收敛，若收敛则求其值：
\SubQ{1}{$\displaystyle\int_{-1}^4\frac{dx}{x^2+x-2}$。}
\SubQ{2}{$\displaystyle\int_{-1}^1x^{-1/3}\,dx$。}
\end{minipage}\par\vfill
\noindent\begin{minipage}{\textwidth}\ProblemHeading{习题}{5.80}
设 $m,n$ 为正整数。证明：
\SubQ{1}{$\displaystyle\int_0^{2\pi}\sin(mx)\cos(nx)\,dx=0$。}
\SubQ{2}{$\displaystyle\int_0^{2\pi}\sin(mx)\sin(nx)\,dx=\int_0^{2\pi}\cos(mx)\cos(nx)\,dx=\begin{cases}\pi,&m=n,\\0,&m\ne n.\end{cases}$}
\end{minipage}\par\vfill
\noindent\begin{minipage}{\textwidth}\ProblemHeading{习题}{5.81}
设 $m,n$ 为正整数，$B(m,n)=\int_0^1x^m(1-x)^n\,dx$。证明：
\SubQ{1}{$B(m,n)=B(n,m)$。}
\SubQ{2}{$B(m,n)=\dfrac{m!n!}{(m+n+1)!}$。}
\end{minipage}\par\vfill
\noindent\begin{minipage}{\textwidth}\ProblemHeading{习题}{5.82}
计算下列积分：
\SubQ{1}{$\displaystyle\int_{1/2}^2(1+x-1/x)e^{x+1/x}\,dx$。}
\SubQ{2}{$\displaystyle\int_0^{n\pi}x|\sin x|\,dx$，其中 $n$ 为非负整数。}
\SubQ{3}{设 $f(x)=\displaystyle\int_x^{x+2\pi}\bigl(1+e^{\sin t}-e^{-\sin t}\bigr)\,dt+\frac1{1+x}\int_0^1f(t)\,dt$（$0\le x\le1$），求 $\int_0^1f(x)\,dx$。}
\end{minipage}\par\vfill
\newpage
\noindent\begin{minipage}{\textwidth}\ProblemHeading{习题}{5.83}
证明对于 $n=1,2,\ldots$，有 $\displaystyle\frac1{2n+2}<\int_0^{\pi/4}\tan^n x\,dx<\frac1{2n}$。
\end{minipage}\par\vfill
\noindent\begin{minipage}{\textwidth}\ProblemHeading{习题}{5.84}
设 $f$ 在 $[a,b]$ 上可积，且 $\int_a^b f(x)\,dx>0$。证明存在非退化区间 $[\alpha,\beta]\subset[a,b]$，使其中每一点都满足 $f(x)>0$。
\end{minipage}\par\vfill
\noindent\begin{minipage}{\textwidth}\ProblemHeading{习题}{5.85}
证明下列结论：
\SubQ{1}{处处连续的偶函数有一个奇的原函数。}
\SubQ{2}{处处连续的奇函数的每个原函数都是偶函数。}
\end{minipage}\par\vfill
\noindent\begin{minipage}{\textwidth}\ProblemHeading{习题}{5.86}
构造一个处处连续的周期函数 $f$，使 $f$ 的所有原函数都不是周期函数。
\end{minipage}\par\vfill
\newpage
\noindent\begin{minipage}{\textwidth}\ProblemHeading{习题}{5.87}
设 $a_0/(n+1)+a_1/n+\cdots+a_n=0$。用积分方法证明多项式 $a_0x^n+a_1x^{n-1}+\cdots+a_n$ 在 $(0,1)$ 内至少有一个零点。
\end{minipage}\par\vfill
\noindent\begin{minipage}{\textwidth}\ProblemHeading{习题}{5.88}
设 $f\in C[0,\pi]$，且 $\int_0^\pi f(x)\sin x\,dx=\int_0^\pi f(x)\cos x\,dx=0$。证明在 $(0,\pi)$ 内存在两个不同的点 $x_1,x_2$，使 $f(x_1)=f(x_2)=0$。
\end{minipage}\par\vfill
\noindent\begin{minipage}{\textwidth}\ProblemHeading{习题}{5.89}
设 $f$ 处处连续，$f(0)=0$，且 $f^{\prime}(0)$ 存在。记 $F(x)=\int_0^1 f(xu)\,du$。证明 $F$ 处处可导，并求 $F^{\prime}(x)$。
\end{minipage}\par\vfill
\noindent\begin{minipage}{\textwidth}\ProblemHeading{习题}{5.90}
完成下列问题：
\SubQ{1}{设 $f(x)=\begin{cases}e^{-x^2},&|x|\le1,\\1,&|x|>1,\end{cases}$，$F(x)=\int_0^x f(t)\,dt$。讨论 $F$ 在哪些点可导。}
\SubQ{2}{设 $f(x)=\int_0^x\cos(1/t)\,dt$（$x\ge0$，在 $t=0$ 按反常积分解释），证明 $f^{\prime}_+(0)=0$。}
\end{minipage}\par\vfill
\newpage
\noindent\begin{minipage}{\textwidth}\ProblemHeading{习题}{5.91}
设 $f$ 处处连续。证明\[\lim_{h\to0}\frac1h\int_a^b\bigl(f(x+h)-f(x)\bigr)\,dx=f(b)-f(a).\]
\end{minipage}\par\vfill
\noindent\begin{minipage}{\textwidth}\ProblemHeading{习题}{5.92}
设 $f\in C^1[a,b]$。证明 $\displaystyle\lim_{\lambda\to+\infty}\int_a^b f(x)\sin(\lambda x)\,dx=0$。
\end{minipage}\par\vfill
\noindent\begin{minipage}{\textwidth}\ProblemHeading{习题}{5.93}
证明 $\displaystyle\lim_{x\to+\infty}\frac1x\int_0^x|\sin t|\,dt=\frac2\pi$。
\end{minipage}\par\vfill
\noindent\begin{minipage}{\textwidth}\ProblemHeading{习题}{5.94}
证明 $\displaystyle\lim_{n\to\infty}\int_0^{\pi/2}\sin^n x\,dx=0$。
\end{minipage}\par\vfill
\newpage
\noindent\begin{minipage}{\textwidth}\ProblemHeading{习题}{5.95}
设 $a<b$，$f\in C[a,b]$ 且 $f\ge0$，$M=\max_{[a,b]}f$。证明 $\displaystyle\lim_{n\to\infty}\left(\int_a^b f^n(x)\,dx\right)^{1/n}=M$。
\end{minipage}\par\vfill
\noindent\begin{minipage}{\textwidth}\ProblemHeading{习题}{5.96}
设 $f$ 定义于 $[1,+\infty)$ 且非负，证明：
\SubQ{1}{若 $f$ 递增，则对每个正整数 $n$，$0\le\sum_{k=1}^n f(k)-\int_1^n f(x)\,dx\le f(n)$。}
\SubQ{2}{若 $f$ 递减，则对每个正整数 $n$，$0\le\sum_{k=1}^n f(k)-\int_1^n f(x)\,dx\le f(1)$，且该差在 $n\to\infty$ 时有极限 $\alpha\in[0,f(1)]$。}
\end{minipage}\par\vfill
\noindent\begin{minipage}{\textwidth}\ProblemHeading{习题}{5.97}
设 $a<b$，$f,g\in C[a,b]$。证明 Cauchy 积分不等式\[\left(\int_a^b f(x)g(x)\,dx\right)^2\le\int_a^b f^2(x)\,dx\int_a^b g^2(x)\,dx,\]并证明等号成立当且仅当 $f,g$ 至少有一个恒为零，或存在常数 $\lambda$ 使 $f=\lambda g$。
\end{minipage}\par\vfill
\noindent\begin{minipage}{\textwidth}\ProblemHeading{习题}{5.98}
设 $f\in C^1[0,1]$。证明对任意 $a\in[0,1]$，有 $\displaystyle|f(a)|\le\int_0^1|f(x)|\,dx+\int_0^1|f^{\prime}(x)|\,dx$。
\end{minipage}\par\vfill
\newpage
\noindent\begin{minipage}{\textwidth}\ProblemHeading{习题}{5.99}
证明 $\displaystyle0.944<\int_0^1\frac{\sin x}{x}\,dx<0.947$，在 $x=0$ 处将被积函数连续延拓为1。
\end{minipage}\par\vfill
\noindent\begin{minipage}{\textwidth}\ProblemHeading{习题}{5.100}
设 $f\in C^1[0,1]$，且 $|f^{\prime}(x)|\le M$。证明对任意正整数 $n$，有\[\left|\int_0^1 f(x)\,dx-\frac1n\sum_{k=1}^n f(k/n)\right|\le\frac M{2n}.\]
\end{minipage}\par\vfill
\noindent\begin{minipage}{\textwidth}\ProblemHeading{习题}{5.101}
设 $f:\mathbb R\to(0,+\infty)$ 可微，且 $|f^{\prime}(x)-f^{\prime}(y)|\le|x-y|$ 对任意实数 $x,y$ 成立。证明 $(f^{\prime}(x))^2<2f(x)$ 对所有 $x\in\mathbb R$ 成立。
\end{minipage}\par\vfill
